\documentclass[12pt]{article}
\usepackage{amsfonts,amsthm,amsmath,amssymb,mathtools}
\usepackage{mathrsfs}
\usepackage{enumitem}
\usepackage{tikz-cd}
\usepackage{geometry}
\usepackage{mlmodern}
\usepackage{xcolor}

\geometry{letterpaper, portrait, margin=1in}

\setcounter{section}{0}

\renewcommand{\thesection}{\S\arabic{section}}
\renewcommand{\thesubsection}{\arabic{section}}

\DeclareMathOperator{\lcm}{lcm}
\DeclareMathOperator{\im}{im}

\newtheorem{thm}{Theorem}
\newenvironment{theorem}[1]{
	\renewcommand{\thethm}{#1}
	\begin{thm}
		}{\end{thm}}

\newtheorem{lma}{Lemma}
\newenvironment{lemma}[1]{
	\renewcommand{\thelma}{#1}
	\begin{lma}
		}{\end{lma}}

\theoremstyle{definition}
\newtheorem{ex}{}
\newenvironment{exercise}[1]{
	\renewcommand{\theex}{#1}
	\begin{ex}
		}{\end{ex}}

\title{Homework 4}
\author{Connor Mooneyhan}
\date{September 29, 2025}

\begin{document}

\maketitle

Disclaimer: Generative AI was used on problems 2(a), 3, and 8(a) after having much difficulty coming up with the fundamental approach for the proofs.

\begin{ex}
	Let $G$ be a group, and let $H, K$ be subgroups of $G$.
	Let $HK = \left\lbrace hk : h \in H, k \in K \right\rbrace$.
	Prove that $HK$ is a subgroup of $G$ if and only if $HK = KH$.
	Conclude that if $H$ is normal, then $HK$ is always a subgroup.
	Also show that if $HK$ is a subgroup, it is the smallest subgroup containing $H$ and $K$.
\end{ex}
\begin{proof}
	If $HK$ is a subgroup of $G$, then given $g_1 = hk$ for some $h \in H$ and $k \in K$, \begin{equation*}
		g_1^{-1} = (hk)^{-1} = k^{-1}h^{-1} \in KH.
	\end{equation*}
	Since $g_1^{-1}$ is also in $HK$, the same argument gives us $(g_1^{-1})^{-1} = g_1 \in KH$.
	Likewise, let $g_2 = kh$ ($h \in H, k \in K$).
	Then $g_2^{-1} = h^{-1}k^{-1} \in HK$, and since $HK$ is closed under inverses, $g_2 \in HK$.

	On the other hand, if $HK = KH$, then given $h_1,h_2 \in H$ and $k_1,k_2 \in K$, we have \begin{align*}
		(h_1k_1)(h_2k_2)^{-1} = (h_1k_1)(k_2^{-1}h_2^{-1}) \in HK
	\end{align*}
	since both terms on the right hand side are in $HK = KH$.

	Indeed, if $H$ is normal, then given $h_1 \in H$ and $k \in K$, there exists $h_2 \in H$ such that \begin{align*}
		         & kh_1k^{-1} = h_2              \\
		\implies & kh_1 = h_2k \in HK            \\
		\implies & h_1k^{-1} = k^{-1}h_1 \in KH,
	\end{align*}
	so $HK = KH$ and thus $HK$ is a subgroup.

	Also, for any subgroup $J$ containing $H$ and $K$, given any $h \in H$ and $k \in K$, $h$ and $k$ are also in $J$, so $hk \in J$, and thus $HK \subset J$.
\end{proof}

\begin{ex}\
	\begin{enumerate}[label=(\alph*)]
		\item Let $G$ be a group.
		      Prove that if $G/Z(G)$ is cyclic then $G$ is abelian.
		\item Suppose $G$ is a group of order $pq$ for primes $p$ and $q$.
		      Prove that either $G$ is abelian or has trivial center.
	\end{enumerate}
\end{ex}
\begin{proof}\
	\begin{enumerate}[label=(\alph*)]
		\item Let $a, b \in G$ and let $xH$ be a generator for $G/Z(G)$.
		      Then $a = x^mg$ and $b = x^nh$ for some $m, n \in \mathbb{Z}^{\geq 0}$ and $g, h \in Z(G)$, so \begin{equation*}
			      ab = x^mgx^nh = x^{m+n}gh = x^{n+m}hg = x^nhx^mg = ba.
		      \end{equation*}
		\item By Lagrange's theorem, $|Z(G)| \in \left\lbrace 1, p, q, pq \right\rbrace$.
		      If $|Z(G)| = p$ (resp. $q$), then $|G/Z(G)| = q$ (resp. $p$), so $G/Z(G)$ is cyclic since it has prime order and part (a) shows that $G$ is abelian.
		      If $|Z(G)| = pq$, then $Z(G) = G$, so $G$ is abelian.
		      Otherwise, $|Z(G)| = 1$, so $Z(G)$ is trivial.
	\end{enumerate}
\end{proof}

\begin{ex}
	Let $G$ be a finite group, and suppose that $H$ is a subgroup of index $p$ where $p$ is the smallest prime that divides $|G|$.
	Prove that $H$ is normal in $G$.
\end{ex}
\begin{proof}
	Let $G$ act on $G/H$ (the set of left-cosets of $H$ in $G$) by left translation.
	This is a homomorphism $\varphi : G \to S_p$, and $G/\ker \varphi \cong \im \varphi$.
	Clearly $|H| \subset \ker \varphi$, so $|H| \leq |\ker \varphi|$.

	Now consider that $|\im \varphi| = |\varphi(G)|$ divides $|G|$ and $|\im\varphi|$ divides $|S_p|=p!$.
	Let $q$ be a prime divisor of $|\im\varphi|$.
	Then $q \geq p$ since $p$ is the smallest prime dividing $|G|$, but also $q$ divides $p!$, so $q \leq p$ since $p$ is the largest prime dividing $p!$.
	Thus all the prime factors of $|\im\varphi|$ are $p$, but there is only one factor of $p$ in $p!$, so in fact $|\im\varphi|=p$.

	Tying this all together, we see that \begin{equation*}
		|G|/|H| = [G : H] = p = |\im\varphi| = |G/\ker\varphi| = |G|/|\ker\varphi|,
	\end{equation*}
	so $|H| = |\ker\varphi|$, and since $H \subset \ker\varphi$, $H = \ker\varphi$, which is normal.
\end{proof}

\begin{ex}
	Suppose $\varphi : G \to H$ is an injective group homomorphism.
	Prove that if $K$ is a subgroup of $G$, then the index of $K$ in $G$ is the same as the index of $\varphi(K)$ in $\varphi(G)$.
\end{ex}
\begin{proof}
	Since $\varphi$ is injective, $|G| = |\varphi(G)|$ and $|K| = |\varphi(K)|$, so by Lagrange's theorem, \begin{equation*}
		[G : K] = \frac{|G|}{|K|} = \frac{|\varphi(G)|}{|\varphi(K)|} = [\varphi(G) : \varphi(K)].
	\end{equation*}
\end{proof}

\begin{ex}
	Let $M$ and $N$ be normal subgroups of a group $G$ such that $G = MN$.
	Prove that $G/(M \cap N) \cong (G/M) \times (G/N)$.
\end{ex}
\begin{proof}
	Define $\varphi : G \to (G/M) \times (G/N)$ by $g \mapsto (gM, gN)$.
	Then given $g, h \in G$, \begin{equation*}
		\varphi(gh) = (ghM, ghN) = (gM, gN)(hM, hN) = \varphi(g)\varphi(h),
	\end{equation*}
	so $\varphi$ is a homomorphism.

	Now let $m \in M$ and $n \in N$.
	Then since $N$ is normal, there exists $n' \in N$ such that \begin{equation*}
		mn'm^{-1} = n \implies mn' = nm.
	\end{equation*}
	Then \begin{align*}
		\varphi(mn') & = (mn'M, mn'N) \\
		             & = (nmM, mn'N)  \\
		             & = (nM, mN).
	\end{align*}
	Thus $\varphi$ is surjective since all cosets of $M$ (resp. $N$) have representatives in $N$ (resp. $M$) because $G = MN = NM$.

	Also, clearly $\ker \varphi = M \cap N$.
	By the first isomorphism theorem, then, $G/(M \cap N) \cong (G/M) \times (G/N)$.
\end{proof}

\begin{ex}
	Let $G = S^1$, the unit circle in $\mathbb{C}$ under multiplication, and let $n \geq 2$ be an integer.
	Prove that the map $\varphi : G \to G$ given by $\varphi(x) = x^n$ is a surjective group homomorphism and identify its kernel.
	Conclude that $G$ is isomorphic to a proper quotient of itself.
\end{ex}
\begin{proof}
	Let $g, h \in G$, and write $g = e^{ix}$ and $h = e^{iy}$.
	Then \begin{equation*}
		\varphi(gh) = \varphi(e^{i(x + y)}) = e^{in(x + y)} = e^{inx}e^{iny} = \varphi(g)\varphi(h),
	\end{equation*}
	so $\varphi$ is a homomorphism.
	Also, given any $e^{iz} \in G$, $\varphi(e^{i z/n}) = e^{iz}$, so $\varphi$ is surjective.
	The kernel of $\varphi$ is $\left\lbrace e^{i (2\pi m)/n} \right\rbrace$ where $0 \leq m < n$, which is nontrivial since $n \geq 2$.
	Thus the first isomorphism theorem lets us conclude that $G$ is isomorphic to a proper quotient of itself.
\end{proof}

\begin{ex}
	Let $\mathsf{C}$ be a category.
	An object $I$ of $\mathsf{C}$ is called \textit{initial} if given any object $X$ of $\mathsf{C}$, there is a unique morphism $I \to X$ in $\mathsf{C}$.
	There is a similar notion of \textit{final} object of a category.
	\begin{enumerate}[label=(\alph*)]
		\item Prove that initial objects in a category are unique up to unique isomorphism.
		\item Prove that the trivial group is both initial and final in the category of groups.
		\item Let $G$ be a group and let $N$ be a normal subgroup of $G$.
		      Let $\mathsf{C}$ be the category whose objects are group homomorphisms $\varphi : G \to H$, such that $N \subset \ker \varphi$.
		      A morphism from $\varphi_1 : G \to H_1$ to $\varphi_2 : G \to H_2$ in $\mathsf{C}$ is a group homomorphism $\psi : H_1 \to H_2$ such that $\varphi_2 = \psi\varphi_1$.
		      Show that $\mathsf{C}$ is a category, and explain why the projection map $\pi : G \to G / N$ is initial in $\mathsf{C}$.
	\end{enumerate}
\end{ex}
\begin{proof}\
	\begin{enumerate}[label=(\alph*)]
		\item Let $I_1, I_2$ be initial in a category $\mathsf{C}$.
		      Then there exist unique isomorphisms $\varphi: I_1 \to I_2$ and $\psi: I_2 \to I_1$, but in fact $\varphi^{-1} : I_2 \to I_1$ and $\psi^{-1} : I_1 \to I_2$ are also isomorphisms since the inverse of an isomorphism is an isomorphism, so uniqueness gives us $\varphi = \psi^{-1}$ and $\psi = \varphi^{-1}$, and we conclude that $I_1$ and $I_2$ are isomorphic up to unique isomorphism.
		\item Since all homomorphisms take the identity to the identity, the only option for a homomorphism from the trivial group to a given group $G$ is the map $\varphi$ sending the identity to the identity, and this is in fact a homomorphism since \begin{equation*}
			      \varphi(e \cdot e) = \varphi(e) = e_G = e_G \cdot e_G = \varphi(e)\cdot\varphi(e).
		      \end{equation*}
		      On the other hand, the only possible map from $G$ to the trivial group is the one that sends everything to the identity, i.e. the trivial homomorphism.
		\item The identity morphism criterion is satisfied by the identity homomorphism on each object.
		      Also, given objects $f : G \to H_1$, $g : G \to H_2$, and $h : G \to H_3$ along with morphisms $\varphi : H_1 \to H_2$ and $\psi : H_2 \to H_3$ such that $\varphi f = g$ and $\psi g = h$, we have $(\psi\varphi) f = \psi g = h$, so composition holds, and clearly the identity morphisms act as identities with respect to composition.

		      The statement that the projection map is initial in $\mathsf{C}$ is merely a restatement of the universal property of quotient groups (or quotient maps, if you prefer) in categorical terms.
		      In fact, the fact that it is initial or final in some category is what justifies the name ``universal property.''
	\end{enumerate}
\end{proof}

\begin{ex}
	Let $G$ be a group, let $H$ be a normal subgroup, let $x \in H$, and let $\mathcal{K} = [x]_G$.
	Our goal is to prove that $\mathcal{K}$ is the union of $[G : HC_G(x)]$ many conjugacy classes of $H$, all of equal size.
	\begin{enumerate}[label=(\alph*)]
		\item Let $g \in G$ and $a \in H$.
		      Show that $g \cdot ([a]_H) = \left\lbrace ghg^{-1} : h \in [a]_H \right\rbrace$ is equal to $[b]_H$ for some $b \in H$.
		      Note that this gives a $G$-action on the set of all conjugacy classes of $H$.
		\item Suppose $\mathcal{K}$ is a conjugacy class in $G$, and that $\mathcal{K} \subset H$.
		      Explain why $H$ acts on $\mathcal{K}$ by conjugation, so that $\mathcal{K}$ is a union of $H$-conjugacy classes.
		      Show that if $[a]_H$ is one such $H$-conjugacy class, then $g \cdot ([a]_H)$ is as well.
		\item Show that $G$ acts transitively on the set of $H$-conjugacy classes contained in $\mathcal{K}$. Also, provide an explicit bijection between $[a]_H$ and $[b]_H$ for $a, b \in \mathcal{K}$.
		\item We now compute the stabilizer of an $H$-conjugacy class.
		      To begin, show that if $g \in \mathrm{Stab}_G([x]_H)$, then $gxg^{-1} \in [x]_H$.
		      Use this to conclude that $\mathrm{Stab}_G([x]_H) \subset HC_G(x)$.
		      Then use Problem 1 to conclude the reverse containment.
		\item Use the Orbit-Stabilizer theorem (Problem 8 on Homework 3) to conclude that the number of $H$-conjugacy classes that union to give $\mathcal{K}$ is equal to $[G : HC_G(x)]$.
		\item For an application, take $G = S_n$ and $H = A_n$.
		      Show that every $S_n$-conjugacy class of $S_n$ contained in $A_n$ is either an $A_n$-conjugacy class, or the union of two $A_n$-conjugacy classes of the same size.
		      Prove that in this case, $[x]_{S_n} = [x]_{A_n}$ if and only if $x$ commutes with some odd permutation.
	\end{enumerate}
\end{ex}
\begin{proof}\
	\begin{enumerate}[label=(\alph*)]
		\item Let $h \in H$, so that $hah^{-1} \in [a]_H$.
		      Then $ghah^{-1}g^{-1} = (ghg^{-1})(gag^{-1})(gh^{-1}g^{-1})$, where all three terms are in $H$ by normality of $H$, so $ghah^{-1}g^{-1} \in [gag^{-1}]_H$.
		      On the flip side, any element $hgag^{-1}h^{-1} = (gh_1g^{-1})(gag^{-1})(gh_2g^{-1}) \in [a]_H$ for some $h_1,h_2 \in H$.
		\item $H$ acts on $\mathcal{K}$ by conjugation since given $h \in H$ and $gkg^{-1} \in \mathcal{K}$, $hgkg^{-1}h^{-1} = (hg)k(hg)^{-1} \in \mathcal{K}$.
		      If $[a]_H$ is an $H$-conjugacy class contained in $\mathcal{K}$, then the above shows that $g \cdot ([a]_H) \subset \mathcal{K}$, and given $hah^{-1} \in [a]_H$ and $g \in H$, $ghah^{-1}g^{-1} = (gh)a(gh)^{-1} \in \mathcal{K}$.
		\item It suffices to provide the stated bijection.
		      Given $a, b \in \mathcal{K}$\dots
		\item If $g \in \mathrm{Stab}_G([x]_H)$, then by definition $[x]_H = g[x]_Hg^{-1}$, and in particular $gxg^{-1} \in [x]_H$.
		      Then $gxg^{-1} = hxh^{-1}$ for some $h \in H$\dots

		      On the other hand, $HC_G(x) \subset \mathrm{Stab}([x]_H)$ by problem 1 since $H \subset \mathrm{Stab}_G([x]_H)$ and $C_G(x) \subset \mathrm{Stab}_G([x]_H)$.
		\item Since the action is transitive, \begin{equation*}
			      |\mathcal{K}| = |\mathrm{Orb}_G(x)| = [G : \mathrm{Stab}_G(x)] = [G : HC_G(x)].
		      \end{equation*}
	\end{enumerate}
\end{proof}

\end{document}
